\section{AIME I 1995}
        \begin{enumerate}
        	\item (\textit{AIME I 1995}) Square $S_{1}$ is $1\times 1.$  For $i\ge 1,$ the lengths of the sides of square $S_{i+1}$ are half the lengths of the sides of square $S_{i},$ two adjacent sides of square $S_{i}$ are perpendicular bisectors of two adjacent sides of square $S_{i+1},$ and the other two sides of square $S_{i+1},$ are the perpendicular bisectors of two adjacent sides of square $S_{i+2}.$  The total area enclosed by at least one of $S_{1}, S_{2}, S_{3}, S_{4}, S_{5}$ can be written in the form $m/n,$ where $m$ and $n$ are relatively prime positive integers.  Find $m-n.$

                    \begin{figure}[H]
                        \centering
                        \includegraphics[width=0.5\textwidth]{Problem/AIME_I_1995_1_Problem_Figure_1.png}
                    \end{figure}
                    

 \textbf{Answer: }255



\textbf{Solution}

The sum of the areas of the \href{https://artofproblemsolving.com/wiki/index.php?title=Square}{squares} if they were not interconnected is a \href{https://artofproblemsolving.com/wiki/index.php?title=Geometric\_sequence}{geometric sequence}: 

 
            \begin{enumerate}[label=(\roman*)]
                \item $1^2 + \left(\frac{1}{2}\right)^2 + \left(\frac{1}{4}\right)^2 + \left(\frac{1}{8}\right)^2 + \left(\frac{1}{16}\right)^2$
            \end{enumerate}
            Then subtract the areas of the intersections, which is $\left(\frac{1}{4}\right)^2 + \ldots + \left(\frac{1}{32}\right)^2$:

 
            \begin{enumerate}[label=(\roman*)]
                \item $1^2 + \left(\frac{1}{2}\right)^2 + \left(\frac{1}{4}\right)^2 + \left(\frac{1}{8}\right)^2 + \left(\frac{1}{16}\right)^2 - \left[\left(\frac{1}{4}\right)^2 + \left(\frac{1}{8}\right)^2 + \left(\frac{1}{16}\right)^2 + \left(\frac{1}{32}\right)^2\right]$\item $= 1 + \frac{1}{2}^2 - \frac{1}{32}^2$
            \end{enumerate}
            The majority of the terms cancel, leaving $1 + \frac{1}{4} - \frac{1}{1024}$, which simplifies down to $\frac{1024 + \left(256 - 1\right)}{1024}$. Thus, $m-n =  \boxed{255}$.

 Alternatively, take the area of the first square and add $\,\frac{3}{4}$ of the areas of  the remaining squares. This results in $1+ \frac{3}{4}\left[\left(\frac{1}{2}\right)^2 + \ldots + \left(\frac{1}{16}\right)^2\right]$, which when simplified will produce the same answer.

 


	\item (\textit{AIME I 1995}) Find the last three digits of the product of the positive roots of
$\sqrt{1995}x^{\log_{1995}x}=x^2.$

 \textbf{Answer: }025



\textbf{Solution 1}

Taking the $\log_{1995}$ (\href{https://artofproblemsolving.com/wiki/index.php?title=Logarithm}{logarithm}) of both sides and then moving to one side yields the \href{https://artofproblemsolving.com/wiki/index.php?title=Quadratic\_equation}{quadratic equation} $2(\log_{1995}x)^2 - 4(\log_{1995}x)  + 1 = 0$. Applying the \href{https://artofproblemsolving.com/wiki/index.php?title=Quadratic\_formula}{quadratic formula} yields that $\log_{1995}x = 1 \pm \frac{\sqrt{2}}{2}$. Thus, the product of the two roots (both of which are positive) is $1995^{1+\sqrt{2}/2} \cdot 1995^{1 - \sqrt{2}/2} = 1995^2$, making the solution $(2000-5)^2 \equiv \boxed{025} \pmod{1000}$.

 \begin{verbatim}
Alternatively, use Vieta's formula to find the sum of the roots for solutions of $\log_{1995}x$ which yields the product of the two roots.
\end{verbatim}





\textbf{Solution 2}

Instead of taking $\log_{1995}$, we take $\log_x$ of both sides and simplify:

 $\log_x(\sqrt{1995}x^{\log_{1995}x})=\log_x(x^{2})$
$\log_x\sqrt{1995}+\log_x x^{\log_{1995}x}=2$
$\dfrac{1}{2} \log_x 1995 + \log_{1995} x = 2$

 We know that $\log_x 1995$ and $\log_{1995} x$ are reciprocals, so let $a=\log_{1995} x$.  Then we have $\dfrac{1}{2}\left(\dfrac{1}{a}\right) + a = 2$.  Multiplying by $2a$ and simplifying gives us $2a^2-4a+1=0$, as shown above.

 Because $a=\log_{1995} x$, $x=1995^a$. By the quadratic formula, the two roots of our equation are $a=\frac{2\pm\sqrt2}{2}$. This means our two roots in terms of $x$ are $1995^\frac{2+\sqrt2}{2}$ and $1995^\frac{2-\sqrt2}{2}.$ Multiplying these gives $1995^2$
$1995^2\pmod{1000}\equiv 995^2\pmod{1000}\equiv (-5)^2\pmod{1000}\equiv 25\pmod{1000}$, so our answer is $\boxed{025}$.

 



\textbf{Solution 3}

Let $y=\log_{1995}x$. Rewriting the equation in terms of $y$, we have
 \[\sqrt{1995}\left(1995^y\right)^y=1995^{2y}\]
 \[1995^{y^2+\frac{1}{2}}=1995^{2y}\]
 \[y^2+\frac{1}{2}=2y\]
 \[2y^2-4y+1=0\]
 \[y=\frac{4\pm\sqrt{16-\left(4\right)\left(2\right)\left(1\right)}}{4}=\frac{4\pm\sqrt{8}}{4}=\frac{2\pm\sqrt{8}}{2}={1\pm\sqrt{2}}\]Thus, the product of the positive roots is $\left(1995^{\frac{2+\sqrt{8}}{2}}\right)\left(1995^{\frac{2-\sqrt{8}}{2}}\right)=1995^2=\left(2000-5\right)^2$, so the last three digits are $\boxed{025}$.

 



\textbf{Solution 4 (Basically Solution 3 but more detail)}

Set \( y = \log_{1995}(x) \).

 Then we have

 
 \[\sqrt{1995}x^{y}=x^2\]

 Also, because \( y = \log_{1995}(x) \), we have \( 1995^y = x \), so, our new definition is

 
 \[\sqrt{1995} \cdot 1995^{y^2}=1995^{2y}\]

 Upon which squaring both sides to remove \( \sqrt{1995} \)

 
 \[1995 \cdot 1995^{2y^2}=1995^{4y}\]
 \[1995^{2y^2+1}=1995^{4y}\]
 \[2y^2+1=4y\]
 \[2y^2-4y+1=0\]Solving the equation yields \( y_1 = \frac{2 + \sqrt{2}}{2} \) and \( y_2 = \frac{2 - \sqrt{2}}{2} \), in which now we want to find \( 1995^{y_1} \cdot 1995^{y_2} \mod 1000 \), or \( 1995^{\frac{2+ \sqrt{2}}{2} + \frac{2 - \sqrt{2}}{2}} \mod 1000 \), or \( 1995^2 \mod 1000 \), or \( (2000-5)^2 \mod 1000 \), or \( (2000^2-10 \cdot 2000+25) \mod 1000 \) or \( (0-0+25) \mod 1000 \) or simply $\boxed{025}$.

 


	\item (\textit{AIME I 1995}) Starting at $(0,0),$ an object moves in the coordinate plane via a sequence of steps, each of length one.  Each step is left, right, up, or down, all four equally likely.  Let $p$ be the probability that the object reaches $(2,2)$ in six or fewer steps.  Given that $p$ can be written in the form $m/n,$ where $m$ and $n$ are relatively prime positive integers, find $m+n.$

 \textbf{Answer: }067



\textbf{Solution 1}

It takes an even number of steps for the object to reach $(2,2)$, so the number of steps the object may have taken is either $4$ or $6$.

 If the object took $4$ steps, then it must have gone two steps <tt>N</tt> and two steps <tt>E</tt>, in some permutation. There are $\frac{4!}{2!2!} = 6$ ways for these four steps of occuring, and the probability is $\frac{6}{4^{4}}$.

 If the object took $6$ steps, then it must have gone two steps <tt>N</tt> and two steps <tt>E</tt>, and an additional pair of moves that would cancel out, either <tt>N/S</tt> or <tt>W/E</tt>. The sequences <tt>N,N,N,E,E,S</tt> can be permuted in $\frac{6!}{3!2!1!} = 60$ ways. However, if the first four steps of the sequence are <tt>N,N,E,E</tt> in some permutation, it would have already reached the point $(2,2)$ in four moves. There are $\frac{4!}{2!2!}$ ways to order those four steps and $2!$ ways to determine the order of the remaining two steps, for a total of $12$ sequences that we have to exclude. This gives $60-12=48$ sequences of steps. There are the same number of sequences for the steps <tt>N,N,E,E,E,W</tt>, so the probability here is $\frac{2 \times 48}{4^6}$.

 The total probability is $\frac{6}{4^4} + \frac{96}{4^6} = \frac{3}{64}$, and $m+n= \boxed{067}$.

 



\textbf{Solution 2}

Let's let the object wander for 6 steps so we get a constant denominator of $4^{6}$

 In the first case, we count how many ways the object can end at (2,2), at the end of 6 steps. We will also count it even if we go to (2,2), and go back to (2,2). So, there are 2 different paths for the object to end at (2,2): 

 1.To go a permutation of R,R,R,U,U,L or 

 2.To go a permutation of R,R,U,U,U,D. 

 There are 60 ways to permute for each case, giving a total of 120 ways for the object to succeed and end at (2,2). In these 120 ways the object could reach (2,2) first and then come back to (2,2). This will be a factor in our second case.

 In the second case, the object can get to (2,2) first in 4 moves, then move away from (2,2) with the remaing 2 moves. So, there are 6 ways to get to (2,2) in 4 moves, then there are 16 ways the object can "move around", but 4 of the ways will return the object back to (2,2). Those 4 ways were already counted in the first case, so we should only count 12 of the 16 ways to prevent over-counting. Thus, there are $12*6 =$ 72 ways in the second case.

 So, in all, there are 120+72 ways for the object to achieve it's goal of moving to (2,2). Put that over our denominator, we get $\frac{192}{4^{6}} = \frac{3}{4^{3}} = \frac{3}{64}$, in which adding the numerator and denominator get us an answer of $\boxed{067}$

 
- AlexLikeMath

 


	\item (\textit{AIME I 1995}) Circles of radius $3$ and $6$ are externally tangent to each other and are internally tangent to a circle of radius $9$. The circle of radius $9$ has a chord that is a common external tangent of the other two circles. Find the square of the length of this chord.


 
 \textbf{Answer: }224



\textbf{Solution 1}

We label the points as following: the centers of the circles of radii $3,6,9$ are $O_3,O_6,O_9$ respectively, and the endpoints of the chord are $P,Q$. Let $A_3,A_6,A_9$ be the feet of the \href{https://artofproblemsolving.com/wiki/index.php?title=Perpendicular}{perpendiculars} from $O_3,O_6,O_9$ to $\overline{PQ}$ (so $A_3,A_6$ are the points of \href{https://artofproblemsolving.com/wiki/index.php?title=Tangent\_(geometry)}{tangency}). Then we note that $\overline{O_3A_3} \parallel \overline{O_6A_6} \parallel \overline{O_9A_9}$, and $O_6O_9 : O_9O_3 = 3:6 = 1:2$. Thus, $O_9A_9 = \frac{2 \cdot O_6A_6 + 1 \cdot O_3A_3}{3} = 5$ (consider similar triangles). Applying the \href{https://artofproblemsolving.com/wiki/index.php?title=Pythagorean\_Theorem}{Pythagorean Theorem} to $\triangle O_9A_9P$, we find that 
 \[PQ^2 = 4(A_9P)^2 = 4[(O_9P)^2-(O_9A_9)^2] = 4[9^2-5^2] = \boxed{224}\]

                    \begin{figure}[H]
                        \centering
                        \includegraphics[width=0.5\textwidth]{Solution/AIME_I_1995_4_Solution_Figure_1.png}
                    \end{figure}
                    




\textbf{Solution 2 (Analytic Geometry)}


                    \begin{figure}[H]
                        \centering
                        \includegraphics[width=0.5\textwidth]{Solution/AIME_I_1995_4_Solution_Figure_2.png}
                    \end{figure}
                    

 Let $A$ be defined as the origin of a coordinate plane with the $y$-axis running across the chord and $C(6\sqrt{2},0)$ by the \href{https://artofproblemsolving.com/wiki/index.php?title=Pythagorean\_Theorem}{Pythagorean Theorem}. Then we have $D(0,-6)$ and $F(6\sqrt{2},-3)$, and since $\frac{DE}{DF}=\frac{1}{3}$, the point $E$ is one-third of the way from $D$ to $F$, so point $E$ has coordinates $(2\sqrt{2},-5)$. $E$ is the center of the circle with radius $9$, so the equation of this circle is $(x-2\sqrt{2})^2+(y+5)^2=81$. Since the chord's equation is $y=0$, we must find all values of $x$ satisfying the equation of the circle such that $y=0$. We find that $x-2\sqrt{2}=\pm\sqrt{56}$, so the chord has length $|\sqrt{56}+2\sqrt{2}-(-\sqrt{56}+2\sqrt{2})|=2\sqrt{56}$ and the answer is $(2\sqrt{56})^2=\boxed{224}$.

 \textbf{Note} For more practice, try 2005 AIME II \#8. The idea is identical and it is just more work. 

 


	\item (\textit{AIME I 1995}) For certain real values of $a, b, c,$ and $d_{},$  the equation $x^4+ax^3+bx^2+cx+d=0$ has four non-real roots.  The product of two of these roots is $13+i$ and the sum of the other two roots is $3+4i,$ where $i=\sqrt{-1}.$  Find $b.$

 \textbf{Answer: }051



\textbf{Solution 1}

Since the \href{https://artofproblemsolving.com/wiki/index.php?title=Coefficient}{coefficients} of the \href{https://artofproblemsolving.com/wiki/index.php?title=Polynomial}{polynomial} are real, it follows that the non-real roots must come in \href{https://artofproblemsolving.com/wiki/index.php?title=Complex\_conjugate}{complex conjugate} pairs. Let the first two roots be $m,n$. Since $m+n$ is not real, $m,n$ are not conjugates, so the other pair of roots must be the conjugates of $m,n$. Let $m'$ be the conjugate of $m$, and $n'$ be the conjugate of $n$. Then,
 \[m\cdot n = 13 + i,m' + n' = 3 + 4i\Longrightarrow m'\cdot n' = 13 - i,m + n = 3 - 4i.\]By \href{https://artofproblemsolving.com/wiki/index.php?title=Vieta\%27s\_formulas}{Vieta's formulas}, we have that $b = mm' + nn' + mn' + nm' + mn + m'n' = (m + n)(m' + n') + mn + m'n' = \boxed{051}$.

 



\textbf{Solution 2}

Let's assume that the 2 roots multiplied together are p+qi, and r+si, and the two roots added together are the conjugates of the previous roots. Using Vieta, we get 
 \[b = (p+qi)(r+si) + (p+qi)(r-si) + (p-qi)(r+si) + (p-qi)(r-si) + (p+qi)(p-qi) + (r+si)(r-si) =\]



	\item (\textit{AIME I 1995}) Let $n=2^{31}3^{19}.$  How many positive integer divisors of $n^2$ are less than $n_{}$ but do not divide $n_{}$?


 
 \textbf{Answer: }589



\textbf{Solution 1}

We know that $n^2 = 2^{62}3^{38}$ must have $(62+1)\times (38+1)$ \href{https://artofproblemsolving.com/wiki/index.php?title=Factor}{factors} by its \href{https://artofproblemsolving.com/wiki/index.php?title=Prime\_factorization}{prime factorization}. If we group all of these factors (excluding $n$) into pairs that multiply to $n^2$, then one factor per pair is less than $n$, and so there are $\frac{63\times 39-1}{2} = 1228$ factors of $n^2$ that are less than $n$. There are $32\times20-1 = 639$ factors of $n$, which clearly are less than $n$, but are still factors of $n$. Therefore, using complementary counting, there are $1228-639=\boxed{589}$ factors of $n^2$ that do not divide $n$.

 



\textbf{Solution 2}

Let $n=p_1^{k_1}p_2^{k_2}$ for some primes $p_1,p_2$. Then $n^2$ has $\frac{(2k_1+1)(2k_2+1)-1}{2}$ factors less than $n$. 

 This simplifies to $\frac{4k_1k_2+2k_1+2k_2}{2}=2k_1k_2+k_1+k_2$.

 The number of factors of $n$ less than $n$ is equal to $(k_1+1)(k_2+1)-1=k_1k_2+k_1+k_2$.

 Thus, our general formula for $n=p_1^{k_1}p_2^{k_2}$ is

 Number of factors that satisfy the above $=(2k_1k_2+k_1+k_2)-(k_1k_2+k_1+k_2)=k_1k_2$

 Incorporating this into our problem gives $19\times31=\boxed{589}$.

 



\textbf{Solution 3}

Consider divisors of $n^2: a,b$ such that $ab=n^2$.WLOG, let $b\ge{a}$ and $b=\frac{n}{a}$

 Then, it is easy to see that $a$ will always be less than $b$ as we go down the divisor list of $n^2$ until we hit $n$. 

 Therefore, the median divisor of $n^2$ is $n$.

 Then, there are $(63)(39)=2457$ divisors of $n^2$. Exactly $\frac{2457-1}{2}=1228$ of these divisors are $<n$

 There are $(32)(20)-1=639$ divisors of $n$ that are $<n$.

 Therefore, the answer is $1228-639=\boxed{589}$.

 


	\item (\textit{AIME I 1995}) Given that $(1+\sin t)(1+\cos t)=5/4$ and


 $(1-\sin t)(1-\cos t)=\frac mn-\sqrt{k},$

 where $k, m,$ and $n_{}$ are positive integers with $m_{}$ and $n_{}$ relatively prime, find $k+m+n.$

 \textbf{Answer: }027



\textbf{Solution 1}

From the givens, $2\sin t \cos t + 2 \sin t + 2 \cos t = \frac{1}{2}$, and adding $\sin^2 t + \cos^2t = 1$ to both sides gives $(\sin t + \cos t)^2 + 2(\sin t + \cos t) = \frac{3}{2}$.  Completing the square on the left in the variable $(\sin t + \cos t)$ gives $\sin t + \cos t = -1 \pm \sqrt{\frac{5}{2}}$.  Since $|\sin t + \cos t| \leq \sqrt 2 < 1 + \sqrt{\frac{5}{2}}$, we have $\sin t + \cos t = \sqrt{\frac{5}{2}} - 1$.  Subtracting twice this from our original equation gives $(\sin t - 1)(\cos t - 1) = \sin t \cos t - \sin t - \cos t + 1 = \frac{13}{4} - \sqrt{10}$, so the answer is $13 + 4 + 10 = \boxed{027}$.

 



\textbf{Solution 2}

Let $(1 - \sin t)(1 - \cos t) = x$. Multiplying $x$ with the given equation, $\frac{5x}{4} = (1 - \sin^2 t)(1 - \cos^2 t) = \sin^2 t \cos ^2 t$, and $\frac{\sqrt{5x}}{2} = \sin t \cos t$. Simplifying and rearranging the given equation, $\sin t + \cos t = \frac{5}{4} - (\sin^2 t + \cos^2 t) - \sin t \cos t = \frac{1}{4} - \frac{\sqrt{5x}}{2}$. Notice that $(1 + \sin t)(1 + \cos t) - 2(\sin t + \cos t) = x$, and substituting, $x = \frac{5}{4} - 2( \frac{1}{4} - \frac{\sqrt{5x}}{2}) = \frac{3}{4} + \sqrt{5x}$. Rearranging and squaring, $5x = x^2 - \frac{3}{2} x + \frac{9}{16}$, so $x^2 - \frac{13}{2} x + \frac{9}{16} = 0$, and $x = \frac{13}{4} \pm \sqrt{10}$, but clearly, $0 \leq x < 4$. Therefore, $x = \frac{13}{4} - \sqrt{10}$, and the answer is $13 + 4 + 10 = \boxed{027}$.

 



\textbf{Solution 3}

We want $1+\sin t \cos t-\sin t-\cos t$. However, note that we only need to find $\sin t+\cos t$.

 Let $y = \sin t+\cos t \rightarrow y^2 = \sin^2 t + \cos^2 t + 2\sin t \cos t = 1 + 2\sin t \cos t$

 From this we have $\sin t \cos t = \frac{y^2-1}{2}$ and $\sin t + \cos t = y$

 Substituting, we have $2y^2+4y-3=0 \rightarrow y = \frac {-2 \pm \sqrt{10}}{2}$
$\frac{5}{4} - 2(\frac{-2+\sqrt{10}}{2}) = \frac{13}{4}-\sqrt{10} \rightarrow 13+10+4=\boxed{027}$.

 


	\item (\textit{AIME I 1995}) For how many ordered pairs of positive integers $(x,y),$ with $y<x\le 100,$ are both $\frac xy$ and $\frac{x+1}{y+1}$ integers?


 
 \textbf{Answer: }085



\textbf{Solution 1}

Since $y|x$, $y+1|x+1$, then $\text{gcd}\,(y,x)=y$ (the bars indicate \href{https://artofproblemsolving.com/wiki/index.php?title=Divisibility}{divisibility}) and $\text{gcd}\,(y+1,x+1)=y+1$. By the \href{https://artofproblemsolving.com/wiki/index.php?title=Euclidean\_algorithm}{Euclidean algorithm}, these can be rewritten respectively as $\text{gcd}\,(y,x-y)=y$ and $\text{gcd}\, (y+1,x-y)=y+1$, which implies that both $y,y+1 | x-y$. Also, as $\text{gcd}\,(y,y+1) = 1$, it follows that $y(y+1)|x-y$. <sup>[\href{https://artofproblemsolving.com\#note\_1}{1}]</sup>

 Thus, for a given value of $y$, we need the number of multiples of $y(y+1)$ from $0$ to $100-y$ (as $x \le 100$, but not including 0 because x cannot be equal to y). It follows that there are $\left\lfloor\frac{100-y}{y(y+1)} \right\rfloor$ satisfactory positive integers for all integers $y \le 100$. The answer is 

 
 \[\sum_{y=1}^{99} \left\lfloor\frac{100-y}{y(y+1)} \right\rfloor = 49 + 16 + 8 + 4 + 3 + 2 + 1 + 1 + 1 = \boxed{085}.\]

 

\href{https://artofproblemsolving.com\#ref\_1}{\textasciicircum{}} Another way of stating this is to note that if $\frac{x}{y}$ and $\frac{x+1}{y+1}$ are integers, then $\frac{x}{y} - 1 = \frac{x-y}{y}$ and $\frac{x+1}{y+1} - 1 = \frac{x-y}{y+1}$ must be integers. Since $y$ and $y+1$ cannot share common prime factors, it follows that $\frac{x-y}{y(y+1)}$ must also be an integer.

 



\textbf{Solution 2}

We know that $x \equiv 0 \mod y$ and  $x+1 \equiv 0 \mod y+1$.

 Write $x$ as $ky$ for some integer $k$. Then, $ky+1 \equiv 0\mod y+1$. We can add $k$ to each side in order to factor out a $y+1$. So, $ky+k+1 \equiv k \mod y+1$ or $k(y+1)+1 \equiv k \mod y+1$. We know that $k(y+1) \equiv 0 \mod y+1$. We finally achieve the congruence $1-k \equiv 0 \mod y+1$. 

 We can now write $k$ as $(y+1)a+1$. Plugging this back in, if we have a value for $y$, then $x = ky = ((y+1)a+1)y = y(y+1)a+y$. We only have to check values of $y$ when $y(y+1)<100$. This yields the equations $x = 2a+1, 6a+2, 12a+3, 20a+4, 30a+5, 42a+6, 56a+7, 72a+8, 90a+9$. 

 Finding all possible values of $a$ such that $y<x<100$, we get $49 + 16 + 8 + 4 + 3 + 2 + 1 + 1 + 1 = \boxed{085}.$




\textbf{Solution 3 (Rigorous, but straightforward)}

We use casework: 

 For $y=1$, we have $3,5,\cdots ,99$, or $49$ cases. 

 For $y=2$, we have $8,14,\cdots ,98$, or $16$ cases. 

 For $y=3$, we have $15,27,\cdots ,99$, or $8$ cases. 

 For $y=4$, we have $24,44\cdots ,84$, or $4$ cases. 

 For $y=5$, we have $35,65,95$, or $3$ cases. 

 For $y=6$, we have $48,90$, or $2$ cases. 

 For $y=7$, we have $63$, or $1$ case. 

 For $y=8$, we have $80$, or $1$ case. 

 For $y=9$, we have $99$, or $1$ case. 

 Adding, we get our final result of $\boxed{085}.$

 Note: In general, all of the solutions for all $y$ are $y+by(y+1)$, for any $b$ such that $y+by(y+1)=x<100$

 Also, note that $n+1$ just starts with $(y+1)^2$, and adds $y(y+1)$ until it is greater than $100$.

 


	\item (\textit{AIME I 1995}) Triangle $ABC$ is isosceles, with $AB=AC$ and altitude $AM=11.$  Suppose that there is a point $D$ on $\overline{AM}$ with $AD=10$ and $\angle BDC=3\angle BAC.$  Then the perimeter of $\triangle ABC$ may be written in the form $a+\sqrt{b},$ where $a$ and $b$ are integers.  Find $a+b.$

                    \begin{figure}[H]
                        \centering
                        \includegraphics[width=0.5\textwidth]{Problem/AIME_I_1995_9_Problem_Figure_1.png}
                    \end{figure}
                    

 \textbf{Answer: }616



\textbf{Solution 1}

Let $x=\angle CAM$, so $3x=\angle CDM$. Then, $\frac{\tan 3x}{\tan x}=\frac{CM/1}{CM/11}=11$. Expanding $\tan 3x$ using the angle sum identity gives 
 \[\tan 3x=\tan(2x+x)=\frac{3\tan x-\tan^3x}{1-3\tan^2x}.\] Thus, $\frac{3-\tan^2x}{1-3\tan^2x}=11$. Solving, we get $\tan x= \frac 12$. Hence, $CM=\frac{11}2$ and $AC= \frac{11\sqrt{5}}2$ by the \href{https://artofproblemsolving.com/wiki/index.php?title=Pythagorean\_Theorem}{Pythagorean Theorem}. The total perimeter is $2(AC + CM) = \sqrt{605}+11$. The answer is thus $a+b=\boxed{616}$.

 



\textbf{Solution 2}

In a similar fashion, we encode the angles as complex numbers, so if $BM=x$, then $\angle BAD=\text{Arg}(11+xi)$ and $\angle BDM=\text{Arg}(1+xi)$. So we need only find $x$ such that $\text{Arg}((11+xi)^3)=\text{Arg}(1331-33x^2+(363x-x^3)i)=\text{Arg}(1+xi)$. This will happen when $\frac{363x-x^3}{1331-33x^2}=x$,  which simplifies to $121x-4x^3=0$. Therefore, $x=\frac{11}{2}$. By the Pythagorean Theorem, $AB=\frac{11\sqrt{5}}{2}$, so the perimeter is $11+11\sqrt{5}=11+\sqrt{605}$, giving us our answer, $\boxed{616}$.

 



\textbf{Solution 3}

Let $\angle BAD=\alpha$, so $\angle BDM=3\alpha$, $\angle BDA=180-3\alpha$, and thus $\angle ABD=2\alpha.$ We can then draw the angle bisector of $\angle ABD$, and let it intersect $\overline{AM}$ at $E.$ Since $\angle BAE=\angle ABE$, $AE=BE.$ Let $AE=x$. Then we see by the Pythagorean Theorem, $BM=\sqrt{BE^2-ME^2}=\sqrt{x^2-(11-x)^2}=\sqrt{22x-121}$, $BD=\sqrt{BM^2+1}=\sqrt{22x-120}$, $BA=\sqrt{BM^2+121}=\sqrt{22x}$, and $DE=10-x.$ By the angle bisector theorem, $BA/BD=EA/ED.$ Substituting in what we know for the lengths of those segments, we see that 
 \[\frac{\sqrt{22x}}{\sqrt{22x-120}}=\frac{x}{10-x}.\] multiplying by both denominators and squaring both sides yields 
 \[22x(10-x)^2=x^2(22x-120)\] which simplifies to $x=\frac{55}{8}.$ Substituting this in for x in the equations for $BA$ and $BM$ yields $BA=\frac{\sqrt{605}}{2}$ and $BM=\frac{11}{2}.$ Thus the perimeter is $11+\sqrt{605}$, and the answer is $\boxed{616}$.

 



\textbf{Solution 4}

The triangle is symmetrical so we can split it in half ($\triangle ABM$ and $\triangle ACM$). 

 Let $\angle BAM = y$ and $\angle BDM = 3y$. By the Law of Sines on triangle $BAD$, $\frac{10}{\sin 2y} = \frac{BD}{\sin y}$. Using $\sin 2y = 2\sin y\cos y$ we can get $BD = \frac{5}{\cos y}$. We can use this information to relate $BD$ to $DM$ by using the Law of Sines on triangle $BMD$.

 
 \[\frac{\frac{5}{\cos y}}{\sin BMD} = \frac{1}{\sin 90^\circ - 3y}\]
$\sin BMD = 1$ (as $\angle BMD$ is a right angle), so $\frac{1}{\sin 90^\circ - 3y} = \frac{5}{\cos y}$. Using the identity $\sin 90^\circ - x = \cos x$, we can turn the equation into::

 
 \[\frac{1}{\cos 3y} = \frac{5}{\cos y}\]

 \[5\cos 3y = \cos y\]

 \[5(4\cos ^3 y - 3\cos y) = \cos y\]

 \[20\cos ^3 y = 16 \cos y\]

 \[5\cos ^3 y = 4\cos y\]

 \[5\cos ^2 y = 4\]

 \[\cos ^2 y = \frac{4}{5}\]

 Now that we've found $\cos y$, we can look at the side lengths of $BM$ and $AB$ (since they are symmetrical, the perimeter of $\triangle ABC$ is $2(BM + AB)$.

 We note that $BM = 11\tan y$ and $AB = 11\sec y$.

 
 \[\sin ^2 y = 1 - \cos ^2 y\]

 \[\sin ^2 y = \frac{1}{5}\]

 \[\tan ^2 y = \frac{1}{4}\]

 \[\tan y = \frac{1}{2}\]

 (Note it is positive since $BM > 0$).

 
 \[\sec ^2 y = \frac{5}{4}\]

 \[\sec y = \frac{\sqrt{5}}{2}\]

 \[BM + AB = 11\frac{\sqrt{5}+1}{2}\]

 \[2(BM + AB) = 11(\sqrt{5} + 1)\]

 \[2(BM + AB) = 11\sqrt{5} + 11\]

 \[2(BM + AB) = \sqrt{605} + 11\]

 The answer is $\boxed{616}$.

 



\textbf{Solution 5}

Suppose $\angle BAM=\angle CAM =x$, since $\angle BDC=3\angle BAC$, we have $\angle BDM=\angle MDC = 3x$. Therefore, $\angle DBC=\angle DCB = 90^\circ -3x$ and $\angle ABD=\angle DCA=2x$. As a result, $\triangle KAC$ is isosceles, $KC=KA$.

 Let $H$ be a point on the extension of $CD$ through $D$ such that $\overline{HB}\perp\overline{BC}$ and denote the intersection of $\overline{HC}$ and $\overline{AB}$ as $K$. Then, $BH=2DM=2, \overline{HB}\parallel\overline{DM}$, and $HD=DC$ by the \href{https://artofproblemsolving.com/wiki/index.php?title=Midpoint}{Midpoint} Theorem. So, $\angle HBA=x$ and $\angle CDM=\angle CHB=\angle HDA= 3x$. 

 Consequently, $\triangle HBK\sim \triangle DAK$, 
 \[\frac{BK}{KA}=\frac{HK}{KD}=\frac{1}{5}\] Assume $BK=a$ and $HK=b$, then $KA=5a$ and $KD = 5b$. Since $KC=KA, KC=5a$, and since $HD=DC$, $KC=11b$. Therefore, $a=\frac{11}{5}b$.

 In $\triangle BDM$, by the \href{https://artofproblemsolving.com/wiki/index.php?title=Pythagorean\_Theorem}{Pythagorean Theorem}, $BM=\sqrt{36b^2-1}$. Similarly in $\triangle BAM$, $BM=\sqrt{36a^2-121}$. So 
 \[\sqrt{36a^2-121}=\sqrt{36b^2-1}\] Since $a=\frac{11}{5}b$, we have $b=\frac{5\sqrt{5}}{12}$ and $a=\frac{11\sqrt{5}}{12}$. Consequently, $BM=\frac{11}{2}$ and $AB=\frac{11\sqrt{5}}{2}$. Thus, the perimeter of $\triangle ABC$ is $11+\sqrt{605}$, and the answer is $\boxed{616}$.

 


	\item (\textit{AIME I 1995}) What is the largest positive integer that is not the sum of a positive integral multiple of 42 and a positive composite integer?


 
 \textbf{Answer: }215



\textbf{Solution}

The requested number $\mod {42}$ must be a prime number. Also, every number that is a multiple of $42$ greater than that prime number must also be prime, except for the requested number itself. So we make a table, listing all the primes up to $42$ and the numbers that are multiples of $42$ greater than them, until they reach a composite number.

 
 \[\begin{tabular}{|r||r|r|r|r|r|} \hline 1 & 43 & 85&&& \\ 2&44&&&& \\ 3&45&&&& \\ 5&47&89&131&173&215 \\ 7&49&&&& \\ 11&53&95&&& \\ 13&55&&&& \\ 17&59&101&143&& \\ 19&61&103&145&& \\ 23&65&&&& \\ 29&71&113&155&& \\ 31&73&115&&& \\ 37&79&121&&& \\  41&83&125&&& \\ \hline \end{tabular}\]

 Since $\boxed{215}$ is the greatest number in the list, it is the answer. Note that considering $\mod {5}$ would have shortened the search, since $\text{gcd}(5,42)=1$, and so within $5$ numbers at least one must be divisible by $5$.

 \textbf{Afterword} Basically, we are looking for a number where when a multiple of 42 is subtracted from it, the result is a prime number. Any number that ends in a 5 is not prime, except for 5 itself. Since 42 keeps the parity the same and advances odd numbers in the unit digit, then we can conclude that the sought number is $5 \mod 42$. Specifically, $5 * 42 + 5$.

 -jackshi2006

 



\textbf{Solution 2}

Let our answer be $n$. Write $n = 42a + b$, where $a, b$ are positive integers and $0 \leq b < 42$. Then note that $b, b + 42, ... , b + 42(a-1)$ are all primes. 

 If $b$ is $0\mod{5}$, then $b = 5$ because $5$ is the only prime divisible by $5$. We get $n = 215$ as our largest possibility in this case. 

 If $b$ is $1\mod{5}$, then $b + 2 \times 42$ is divisible by $5$ and thus $a \leq 2$. Thus, $n \leq 3 \times 42 = 126 < 215$.

 If $b$ is $2\mod{5}$, then $b + 4 \times 42$ is divisible by $5$ and thus $a \leq 4$. Thus, $n \leq 5 \times 42 = 210 < 215$.

 If $b$ is $3\mod{5}$, then $b + 1 \times 42$ is divisible by $5$ and thus $a = 1$. Thus, $n \leq 2 \times 42 = 84 < 215$.

 If $b$ is $4\mod{5}$, then $b + 3 \times 42$ is divisible by $5$ and thus $a \leq 3$. Thus, $n \leq 4 \times 42 = 168 < 215$.

 Our answer is $\boxed{215}$.

 


	\item (\textit{AIME I 1995}) A right rectangular prism $P_{}$ (i.e., a rectangular parallelepiped) has sides of integral length $a, b, c,$ with $a\le b\le c.$  A plane parallel to one of the faces of $P_{}$ cuts $P_{}$ into two prisms, one of which is similar to $P_{},$ and both of which have nonzero volume.  Given that $b=1995,$ for how many ordered triples $(a, b, c)$ does such a plane exist?


 
 \textbf{Answer: }040



\textbf{Solution}

Let $P'$ be the prism similar to $P$, and let the sides of $P'$ be of length $x,y,z$, such that $x \le y \le z$. Then 

 
 \[\frac{x}{a} = \frac{y}{b} = \frac zc < 1.\] 

 Note that if the ratio of similarity was equal to $1$, we would have a prism with zero volume. As one face of $P'$ is a face of $P$, it follows that $P$ and $P'$ share at least two side lengths in common. Since $x < a, y < b, z < c$, it follows that the only possibility is $y=a,z=b=1995$. Then,

 
 \[\frac{x}{a} = \frac{a}{1995} = \frac{1995}{c} \Longrightarrow ac = 1995^2 = 3^25^27^219^2.\] 

 The number of factors of $3^25^27^219^2$ is $(2+1)(2+1)(2+1)(2+1) = 81$. Only in $\left\lfloor \frac {81}2 \right\rfloor = 40$ of these cases is $a < c$ (for $a=c$, we end with a prism of zero volume). We can easily verify that these will yield nondegenerate prisms, so the answer is $\boxed{040}$.

 


	\item (\textit{AIME I 1995}) Pyramid $OABCD$ has square base $ABCD,$ congruent edges $\overline{OA}, \overline{OB}, \overline{OC},$ and $\overline{OD},$ and $\angle AOB=45^\circ.$  Let $\theta$ be the measure of the dihedral angle formed by faces $OAB$ and $OBC.$  Given that $\cos \theta=m+\sqrt{n},$ where $m_{}$ and $n_{}$ are integers, find $m+n.$

 \textbf{Answer: }005



\textbf{Solution 1 (trigonometry)}


                    \begin{figure}[H]
                        \centering
                        \includegraphics[width=0.5\textwidth]{Solution/AIME_I_1995_12_Solution_Figure_1.png}
                    \end{figure}
                    

 The angle $\theta$ is the angle formed by two \href{https://artofproblemsolving.com/wiki/index.php?title=Perpendicular}{perpendiculars} drawn to 

 


	\item (\textit{AIME I 1995}) Let $f(n)$ be the integer closest to $\sqrt[4]{n}.$  Find $\sum_{k=1}^{1995}\frac 1{f(k)}.$

 \textbf{Answer: }400



\textbf{Solution}

When $\left(k - \frac {1}{2}\right)^4 \leq n < \left(k + \frac {1}{2}\right)^4$, $f(n) = k$. Thus there are $\left \lfloor \left(k + \frac {1}{2}\right)^4 - \left(k - \frac {1}{2}\right)^4 \right\rfloor$ values of $n$ for which $f(n) = k$. Expanding using the \href{https://artofproblemsolving.com/wiki/index.php?title=Binomial\_theorem}{binomial theorem}, 

 
 \begin{align*} \left(k + \frac {1}{2}\right)^4 - \left(k - \frac {1}{2}\right)^4 &= \left(k^4 + 2k^3 + \frac 32k^2 + \frac 12k + \frac 1{16}\right) - \left(k^4 - 2k^3 + \frac 32k^2 - \frac 12k + \frac 1{16}\right)\\ &= 4k^3 + k. \end{align*}
 

 Thus, $\frac{1}{k}$ appears in the summation $4k^3 + k$ times, and the sum for each $k$ is then $(4k^3 + k) \cdot \frac{1}{k} = 4k^2 + 1$. From $k = 1$ to $k = 6$, we get $\sum_{k=1}^{6} 4k^2 + 1 = 364 + 6 = 370$ (either adding or using the \href{https://artofproblemsolving.com/wiki/index.php?title=Perfect\_square}{sum of consecutive squares formula}).  

 But this only accounts for $\sum_{k = 1}^{6} (4k^3 + k) = 4\left(\frac{6(6+1)}{2}\right)^2 + \frac{6(6+1)}{2} = 1764 + 21 = 1785$ terms, so we still have $1995 - 1785 = 210$ terms with $f(n) = 7$. This adds $210 \cdot \frac {1}{7} = 30$ to our summation, giving $\boxed{400}$.

 



\textbf{Solution 2}

This is a pretty easy problem just to bash. Since the max number we can get is $7$, we just need to test $n$ values for $1.5,2.5,3.5,4.5,5.5$ and $6.5$. Then just do how many numbers there are times $\frac{1}{\lfloor n \rfloor}$, which should be $5+17+37+65+101+145+30 = \boxed{400}$



	\item (\textit{AIME I 1995}) In a circle of radius 42, two chords of length 78 intersect at a point whose distance from the center is 18.  The two chords divide the interior of the circle into four regions.  Two of these regions are bordered by segments of unequal lengths, and the area of either of them can be expressed uniquely in the form $m\pi-n\sqrt{d},$ where $m, n,$ and $d_{}$ are positive integers and $d_{}$ is not divisible by the square of any prime number.  Find $m+n+d.$

 \textbf{Answer: }378



\textbf{Solution}

Let the center of the circle be $O$, and the two chords be $\overline{AB}, \overline{CD}$ and intersecting at $E$, such that $AE = CE < BE = DE$. Let $F$ be the midpoint of $\overline{AB}$. Then $\overline{OF} \perp \overline{AB}$. 

 
                    \begin{figure}[H]
                        \centering
                        \includegraphics[width=0.5\textwidth]{Solution/AIME_I_1995_14_Solution_Figure_1.png}
                    \end{figure}
                    

 By the \href{https://artofproblemsolving.com/wiki/index.php?title=Pythagorean\_Theorem}{Pythagorean Theorem}, $OF = \sqrt{OB^2 - BF^2} = \sqrt{42^2 - 39^2} = 9\sqrt{3}$, and $EF = \sqrt{OE^2 - OF^2} = 9$. Then $OEF$ is a $30-60-90$ \href{https://artofproblemsolving.com/wiki/index.php?title=Right\_triangle}{right triangle}, so $\angle OEB = \angle OED = 60^{\circ}$. Thus $\angle BEC = 60^{\circ}$, and by the \href{https://artofproblemsolving.com/wiki/index.php?title=Law\_of\_Cosines}{Law of Cosines},

 $BC^2 = BE^2 + CE^2 - 2 \cdot BE \cdot CE \cos 60^{\circ} = 42^2.$

 It follows that $\triangle BCO$ is an \href{https://artofproblemsolving.com/wiki/index.php?title=Equilateral\_triangle}{equilateral triangle}, so $\angle BOC = 60^{\circ}$. The desired area can be broken up into two regions, $\triangle BCE$ and the region bounded by $\overline{BC}$ and minor arc $\stackrel{\frown}{BC}$. The former can be found by \href{https://artofproblemsolving.com/wiki/index.php?title=Heron\%27s\_formula}{Heron's formula} to be $[BCE] = \sqrt{60(60-48)(60-42)(60-30)} = 360\sqrt{3}$. The latter is the difference between the area of \href{https://artofproblemsolving.com/wiki/index.php?title=Sector}{sector} $BOC$ and the equilateral $\triangle BOC$, or $\frac{1}{6}\pi (42)^2 - \frac{42^2 \sqrt{3}}{4} = 294\pi - 441\sqrt{3}$. 

 Thus, the desired area is $360\sqrt{3} + 294\pi - 441\sqrt{3} = 294\pi - 81\sqrt{3}$, and $m+n+d = \boxed{378}$.

 Note: the area of $\triangle BCE$ can be more easily found by using the sine method $[\triangle] = \frac{1}{2} ab \sin C$. $[BCE] = 30 \cdot 48 \cdot \frac{1}{2} \cdot \sin 60^\circ = 30 \cdot 24 \cdot \frac{\sqrt{3}}{2} = 360\sqrt{3}$

 -NL008

 


	\item (\textit{AIME I 1995}) Let $p_{}$ be the probability that, in the process of repeatedly flipping a fair coin, one will encounter a run of 5 heads before one encounters a run of 2 tails.  Given that $p_{}$ can be written in the form $m/n$ where $m_{}$ and $n_{}$ are relatively prime positive integers, find $m+n$.


 
 \textbf{Answer: }037



\textbf{Solution 1}

Think of the problem as a sequence of <tt>H</tt>'s and <tt>T</tt>'s. No two <tt>T</tt>'s can occur in a row, so the successful sequences are composed of blocks of $1$ to $4$ <tt>H</tt>'s separated by <tt>T</tt>'s and end with $5$ <tt>H</tt>'s. Since the probability that the sequence starts with <tt>TH</tt> is $1/4$, the total probability is that $3/2$ of the probability given that the sequence starts with an <tt>H</tt>. 

 The answer to the problem is then the sum of all numbers of the form $\frac 32 \left( \frac 1{2^a} \cdot \frac 12 \cdot \frac 1{2^b} \cdot \frac 12 \cdot \frac 1{2^c} \cdots \right) \cdot \left(\frac 12\right)^5$, where $a,b,c \ldots$ are all numbers $1-4$, since the blocks of <tt>H</tt>'s can range from $1-4$ in length. The sum of all numbers of the form $(1/2)^a$ is $1/2+1/4+1/8+1/16=15/16$, so if there are n blocks of <tt>H</tt>'s before the final five <tt>H</tt>'s, the answer can be rewritten as the sum of all numbers of the form $\frac 32\left( \left(\frac {15}{16}\right)^n \cdot \left(\frac 12\right)^n \right) \cdot \left(\frac 1{32}\right)=\frac 3{64}\left(\frac{15}{32}\right)^n$, where $n$ ranges from $0$ to $\infty$, since that's how many blocks of <tt>H</tt>'s there can be before the final five. This is an infinite geometric series whose sum is $\frac{3/64}{1-(15/32)}=\frac{3}{34}$, so the answer is $\boxed{037}$.

 



\textbf{Solution 2}

Let $p_H, p_T$ respectively denote the probabilities that a string beginning with <tt>H</tt>'s(H-string) and <tt>T</tt>'s(T-string) are successful. Notice that a T-string is just a T followed by an H-string. Thus, 

 $p_T = \frac 12p_H.$

 A successful string can either be the string HHHHH, start with 1 to 4 H's and then revert to a T, start with a T and then continue with a string starting with <tt>H</tt> (as there cannot be $2$ <tt>T</tt>'s in a row).

 There is a $\frac{1}{2^5} = \frac{1}{32}$ probability we roll <tt>HHHHH</tt>. On the other hand, there is a $\frac{15}{16}$ probability we roll a <tt>H, HH, HHH, or HHHH</tt>, and then a probability of $p_T$ to follow that up with a string starting with <tt>T</tt>. Thus, 

 $p_H = \left(\frac{15}{16}\right) \cdot \left(\frac 12\right) p_H + \frac{1}{32} \Longrightarrow p_H = \frac{1}{17}.$

 The answer is $p_H + p_T = \frac{3}{2}p_H = \frac{3}{34}$, and $m+n = \boxed{037}$.

 



\textbf{Solution 3}

For simplicity, let's compute the complement, namely the probability of getting to $2$ tails before $5$ heads.

 Let $h_{i}$ denote the probability that we get $2$ tails before $5$ heads, given that we have $i$ consecutive heads. Similarly, let $t_{i}$ denote the probability that we get $2$ tails before $5$ heads, given that we have $i$ consecutive tails. Specifically, $h_{5} = 0$ and $t_{2} = 1$. If we can solve for $h_{1}$ and $t_{1}$, we are done; the answer is simply $1/2 * (h_{1} + t_{1})$, since on our first roll, we have equal chances of getting a string with "1 consecutive head" or "1 consecutive tail."

 Consider solving for $t_{1}$. If we have 1 consecutive tail, then (a) rolling a head gets us to 1 consecutive head and (b) rolling a tail gets us to 2 consecutive tails. So, we must have:

 $t_{1} = \frac{1}{2} t_{2} + \frac{1}{2} h_{1}$

 

 Applying similar logic, we get the equations:

 
 \begin{align*} h_{1} &= \frac{1}{2} h_{2} + \frac{1}{2} t_{1}\\ h_{2} &= \frac{1}{2} h_{3} + \frac{1}{2} t_{1}\\ h_{3} &= \frac{1}{2} h_{4} + \frac{1}{2} t_{1}\\ h_{4} &= \frac{1}{2} h_{5} + \frac{1}{2} t_{1} \end{align*}


 Since $h_{5} = 0$, we get $h_{4} = \frac{1}{2} t_{1}$ $\Rightarrow h_{3} = \frac{3}{4} t_{1}$ $\Rightarrow h_{2} = \frac{7}{8} t_{1}$ $\Rightarrow h_{1} = \frac{15}{16} t_{1}$ $\Rightarrow t_{1} = \frac{1}{2} t_{2} + \frac{1}{2} \cdot \frac{15}{16} t_{1} = \frac{1}{2} + \frac{15}{32} t_{1} \Rightarrow t_{1} = \frac{16}{17}$ $\Rightarrow h_{1} = \frac{15}{16} \cdot \frac{16}{17} = \frac{15}{17}$.

 So, the probability of reaching $2$ tails before $5$ heads is $\frac{1}{2} (h_{1} + t_{1}) = \frac{31}{34}$; we want the complement, $\frac{3}{34}$, yielding an answer of $3 + 34 = \boxed{037}$.

 Note: The same approach still works if we tried solving the original problem rather than the complement; we would have simply used $h_{5} = 1$ and $t_{2} = 0$ instead. The repeated back-substitution is cleaner because we used the complement.

 



\textbf{Solution 4(fast)}

Consider what happens in the "endgame" or what ultimately leads to the end. Let A denote that a head has been flipped, and let B denote that a tail has been flipped. The endgame outcomes are AAAAA, BAAAAA, BB, ABB, AABB, AAABB, AAAABB. The probabilities of each of these are $\frac{1}{32},\frac{1}{64},\frac{1}{4},\frac{1}{8},\frac{1}{16},\frac{1}{32},\frac{1}{64}$ respectively. The ones where there is five heads are AAAAA and BAAAAA. The sum of these probabilities are $\frac{3}{64}$. The sum of all these endgame outcomes are $\frac{34}{64}$, hence the desired probability is $\frac{3}{34}$, and in this case $m=3,n=34$ so we have $m+n=\boxed{037}$-vsamc

 



\textbf{Solution 5}

There can be 1-4 heads and 1 tails in any grouping, so we write $h_5,th_5,h_nth_5,th_nth_5,\ldots$ etc., where $1\leq n\leq 4$ and subscript $n$ means there are $n$ of that state in a row. We see that this gives us 2 geometric sequences: one with first term $\frac{1}{2^5}$ and common ratio $\frac{1}{2}*\left(1/2+1/4+1/8+1/16\right)$ and one with first term $\frac{1}{2^6}$ and the same common ratio. Therefore we get $\frac{3}{34}$ or $\fbox{037}$~joeythetoey

 


\end{enumerate}